\documentclass[a4paper]{article}

\usepackage[left=2cm, right=2cm, top=2cm, bottom=2cm]{geometry} 
\usepackage[utf8]{inputenc} 
\usepackage{german} 
\usepackage{amsmath} 
\usepackage{amssymb} 
\usepackage{listings} 
\usepackage{hyperref} 
\usepackage{graphicx} 
\usepackage{eurosym} 
\usepackage{ulem} \DeclareUnicodeCharacter{20AC}{\euro}

% \usepackage{helvet}
\renewcommand{\familydefault}{\sfdefault}

\begin{document}
\thispagestyle{empty}
    
\begin{center}
    \begin{minipage}{.5\textwidth}
      \flushleft
      \includegraphics[width=.4\linewidth]{Regionsbadge.pdf}
    \end{minipage}%
    \begin{minipage}{.5\textwidth}
      \flushright
      \includegraphics[width=.4\linewidth]{Zeichen.pdf}
    \end{minipage}
    % ~\\[1cm]
    
    
    
	{\Huge \textbf{Stammesordnung}}\\
    \small{Arbeitsversion von Jonas Höchst, \today}
    ~\\[2cm]
    \includegraphics[width=0.8\textwidth]{Lutherrose.pdf}
    \\[2,5cm] 
    
    \normalsize{Verband Christlicher Pfadfinderinnen und Pfadfinder}\\
	\Huge{\textbf{Stamm Martin Luther Lumdatal}}
\end{center}


\clearpage
\tableofcontents 
\clearpage

\section{Aufgabe und Mitgliedschaft} % (fold)
\label{sec:aufgabe_und_mitgliedschaft}
	\begin{itemize}
		\item Der VCP Stamm Martin Luther Lumdatal ist ein Stamm des Verbandes Christlicher Pfadfinderinnen und Pfadfinder (VCP) mit Sitz in 34221 Kassel, Wichernweg 3. 
		\item Die Aufgaben und Ziele richten sich nach der Bundesordnung des VCP und werden gegebenen falls durch die bestehende Landesordnung des Landes Hessen sowie die Regionsordnung der Region Kurhessen ergänzt. 
		\item Mitglied des VCP Stammes Martin Luther Lumdatal ist, wer beim Bundesverband angemeldetes Mitglied ist und dort als zum Stamm Martin Luther gehörig registriert ist. Es gelten die Regelungen zu Beginn und Ende der Mitgliedschaft im VCP. 
	\end{itemize}
% section aufgabe_und_mitgliedschaft (end)

\section{Allgemeines} % (fold)
\label{sec:allgemeines}
	\begin{itemize}
		\item Der Stamm Martin Luther Lumdatal hat seinen Sitz in der evangelischen Kirchengemeinde Londorf, ist jedoch offen für konfessionell anders- bzw. nichtgebundene Jugendliche. 
		\item Der Stamm Martin Luther Lumdatal hat sich selbst ein Zeichen gegeben. Das Stammeswappen ist nach dem Namensgeber des Stammes die Lutherrose. 
		\item Der VCP ist ein koedukativer Verband, in dem beide Geschlechter gleichberechtigt sind. Für die bessere Verständlichkeit wird in der Stammesordnung allerdings nicht nach Geschlechtern unterschieden.
	\end{itemize}
% section allgemeines (end)

\section{Die Gruppen des Stammes} % (fold)
\label{sec:die_gruppen_des_stammes}
In dieser Ordnung werden nur strukturelle Aspekte der Arbeit erläutert, näheres zu den Arbeitsweisen der einzelnen Stufen ist in der Stufenkonzeption des VCP nachzulesen. Der Stamm Martin Luther Lumdatal erkennt die Stufenkonzeption des VCP ohne Einschränkungen an und versucht sie bestmöglich in der aktiven Arbeit umzusetzen. Daraus ergeben sich folgende Gruppen für den Stamm.
    \subsection{Die Kinderstufe (Meute)} % (fold)
    \label{sub:die_kinderstufe_meute}
		\begin{itemize}
			\item Die Kinderstufe nimmt Kinder im Alter von 7 bis 10 Jahren auf. Mitglieder der Meute werden als Wölflinge bezeichnet. In Ausnahmefällen kann die Aufnahme von Kindern im Alter von unter 7 Jahren durch die Führungsrunde genehmigt werden. 
			\item Die Arbeitsweisen der Kinderstufe sind auf Rudyard Kiplings Roman „Das Dschungelbuch“ gestützt. Daraus ergeben sich auch die Bezeichnungen in dieser Stufe. 
			\item Die gesamte Kinderstufe des Stammes bezeichnet sich als Meute und kann sich gegebenenfalls in mehrere Rudel unterteilen. 
			\item Der Leiter der Meute wird als Akela bezeichnet. Bei seiner Arbeit unterstützen ihn die Meutenhelfer. In den Aufgabenbereich des Akela fallen die Planung und Durchführung der Meutenstunden sowie sämtlicher anderer Aktion der Meute. 
		\end{itemize}
    % subsection die_kinderstufe_meute (end)
    
    \subsection{Die Jungpfadfinderstufe} % (fold)
    \label{sub:die_jungpfadfinderstufe}
    \begin{itemize}
        \item Die Jungpfadfinderstufe nimmt Kinder im Alter von 10 bis 13 Jahren auf. Mitglieder dieser
Stufe werden als Jungpfadfinder bezeichnet.
        \item Die Arbeit erfolgt in kleinen Gruppen von 4 bis 10 Personen, die Sippen genannt werden. Alle
Jungpfadfinder- und Pfadfindersippen des Stammes werden als Trupp bezeichnet.
        \item Jede Sippe besitzt einen Sippenführer, der für die Planung und Durchführung der
Sippenstunden sowie sämtlicher anderer Aktionen der Sippe zuständig ist.
    \end{itemize}
    % subsection die_jungpfadfinderstufe (end)
    
    \subsubsection{Die Pfadfinderstufe} % (fold)
    \label{ssub:die_pfadfinderstufe}
    \begin{itemize}
        \item Die Pfadfinderstufe nimmt Jugendliche im Alter von 13 bis 16 Jahren auf. Mitglieder dieser
Stufe werden als Pfadfinder bezeichnet.
        \item Die Strukturen der Jungpfadfinderstufe bleiben erhalten.
    \end{itemize}
    % subsubsection die_pfadfinderstufe (end)

    \subsection{Ranger/Rover Stufe} % (fold)
    \label{sub:ranger_rover_stufe}
        \begin{itemize}
            \item Die Ranger/Rover Stufe nimmt Jugendliche im Alter von 16-21 Jahren auf. Sie übernimmt überwiegend leitende Aufgaben im Stamm und organisiert sich selbstständig in einer oder mehreren Rover-Runden. 
            \item Die Ranger/Rover erarbeiten sich eigene Projekte und nehmen überwiegend an überregionalen Veranstaltungen des VCP teil. 
            \item Die Rover-Runden wählen je einen Roversprecher aus sich selbst heraus, der die Gruppe in der Führungsrunde vertritt.
        \end{itemize}
    % subsection ranger_rover_stufe (end)
    
    \subsection{Erwachsenenarbeit} % (fold)
    \label{sub:erwachsenenarbeit}
	\begin{itemize}
		\item Alle Mitglieder des Stammes Martin Luther Lumdatal ab 21 Jahren zählen zur Erwachsenenarbeit. 
		\item Die Form der Erwachsenenarbeit ist je nach Interesse zu gestalten. 
	\end{itemize}
    % subsection erwachsenenarbeit (end)
% section die_gruppen_des_stammes (end)

\section{Führungsrunde und Beauftragungen} % (fold)
\label{sec:fuhrungsrunde_und_beauftragungen}
    \subsection{Führungsrunde} % (fold)
    \label{sub:fuhrungsrunde}
	\begin{itemize}
		\item Die Führungsrunde setzt sich aus der Stammesleitung und den aktiven Gruppenleitern zusammen. 
		\item Freie Mitarbeiter und Meutenhelfer können von der Führungsrunde eingeladen werden. 
		\item Die Führungsrunde wird vom Stammesführer einberufen und sollte mindestens alle 2 Monate stattfinden. 
		\item Die Führungsrunde versucht nach Möglichkeit ohne Abstimmung Einigungen zu finden. Sollte es dennoch zu einen Abstimmungsverfahren kommen, so gilt die einfache Mehrheit ohne Berücksichtigung der Enthaltungen als beschlussfähig. 
		\item Die Führungsrunde hat das Recht, Punkte der Stammesordnung außer Kraft zu setzen. Änderungen der Stammesordnung dürfen nur mit Zustimmung der Mitgliederversammlung erfolgen. 
		\item Des weiteren unterliegt die Führungsrunde einer Schweigepflicht und ist somit auch für vertrauliche Gespräche zugänglich. Daraus ergibt sich das Recht der Führungsrunde Personen aus der Debatte zu bestimmten Themen auszuschließen. 
		\item Die Protokolle der Führungsrunde sind nur für Mitglieder der Führungsrunde einsehbar.
        \item Die Führungsrunde kann das Protokoll auf Anfrage zur Einsicht freigeben.
		\subsubsection{Aufgaben der FüRu} % (fold)
		\label{ssub:aufgaben_der_furu}
		\begin{itemize}
		    \item Aufstellen eines Jahresplanes 
			\item Vorbereitung der Stammesaktivitäten 
			\item Öffentlichkeitsarbeit 
			\item innerverbandliche Informationsvermittlung 
			\item Bewilligung von Mitteln aus der Stammeskasse 
			\item Beauftragen von Vertretern für Gremien der Region und des Landes 
		\end{itemize}
		% subsubsection aufgaben_der_furu (end)
	\end{itemize}
    % subsection fuhrungsrunde (end)
    
    \subsection{Stammesleitung} % (fold)
    \label{sub:stammesleitung}
    Die Stammesleitung besteht aus dem Stammesführer, dem Materialwart, dem Koch und dem Kassenwart. Jedes Amt der Stammesleitung wird von der Mitgliederversammlung gewählt.
    \subsubsection{Stammesführer (StaFü)} % (fold)
    \label{ssub:stammesfuhrer_stafu}
	\begin{itemize}
        \item Der Stammesführer muss volljährig und Mitglied im VCP sein.
        \item Die Amtszeit beträgt 2 Jahre. Entlastung und Wiederwahl ist durch die Mitgliederversammlung möglich.
        \item Der Stammesführer hat den Vorstand über die Führungsrunde. Er beruft die Führungsrunden ein und moderiert diese.
        \item Der Stammesführer ist ermächtigt, zwischen den Führungsrunden weitreichende Entscheidungen den Stamm betreffend eigenmächtig zu entscheiden, muss sich für diese aber der Mitgliederversammlung und der Führungsrunde gegenüber verantworten.
        \item Der Stammesführer hat die Pflicht an den Regionsräten teilzunehmen.
        \item Kann der Stammesführer bestimmte Aufgaben nicht wahrnehmen, ist es ihm möglich diese an Mitglieder der Führungsrunde und freie Mitarbeiter zu delegieren.
	\end{itemize}
    % subsubsection stammesfuhrer_stafu (end)
    
    \subsubsection{Materialwart} % (fold)
    \label{ssub:materialwart}
	\begin{itemize}
		\item Der Materialwart muss nicht volljährig, aber Mitglied im VCP sein. 
		\item Der Materialwart verwaltet das Material. 
		\item Die Amtszeit beträgt 2 Jahre. 
		\item Eine Ausgabe von Material für privaten Gebrauch kann nur mit Genehmigung der Führungsrunde erfolgen. Der Abholer haftet mit seiner Unterschrift für das Material. 
		\item Der Materialwart gibt der Mitgliederversammlung jährlich einen Bericht. 
	\end{itemize}
    % subsubsection materialwart (end)
    
    \subsubsection{Kassenwart} % (fold)
    \label{ssub:kassenwart}
	\begin{itemize}
		\item Der Kassenwart muss volljährig, aber kein Mitglied im VCP sein. 
		\item Der Kassenwart verwaltet die Stammeskasse. Er wird von der Mitgliederversammlung gewählt, beauftragt und entlastet. 
		\item Die Amtszeit beträgt ein Jahr. Eine Wiederwahl ist möglich. 
		\item Der Kassenwart muss der Mitgliederversammlung einmal im Jahr Rechenschaft ablegen. 
	\end{itemize}
    % subsubsection kassenwart (end)
    
    \subsubsection{Koch} % (fold)
    \label{ssub:koch}
	\begin{itemize}
		\item Der Koch muss Mitglied im VCP sein. 
		\item Der Koch ist auf Veranstaltungen für das leibliche Wohl zuständig. Er plant die Verpflegung und koordiniert die Küche. 
		\item Die Amtszeit beträgt ein Jahr. Eine Wiederwahl ist möglich. 
	\end{itemize}
    % subsubsection koch (end)
    
    % subsection stammesleitung (end)
    
    \subsection{Kassenprüfer} % (fold)
    \label{sub:kassenprufer}
	\begin{itemize}
		\item Die Kassenprüfer prüfen die Stammes-, die Hand- und falls bestehend die Sippenkassen. Sie werden von der Mitgliederversammlung gewählt. 
		\item Die Kassenprüfer müssen volljährig, aber kein Mitglied im VCP sein. 
		\item Sie berichten der Mitgliederversammlung über Prüfung der Kassenführung, vor der Entlastung des Kassenwartes. 
		\item Die Amtszeit beträgt ein Jahr. Eine Wiederwahl ist möglich. 
	\end{itemize}
    % subsection kassenprufer (end)
    
    \subsection{Mitgliederversammlungsvorstand (MVV)} % (fold)
    \label{sub:mitgliederversammlungsvorstand_mvv}
	\begin{itemize}
		\item Der MVV besteht aus zwei Personen, die volljährig, Mitglied im VCP und für keine andere Aufgabe durch die MV beauftragt sind. 
		\item Die Amtszeit beträgt 2 Jahre, nach Möglichkeit sollte jedes Jahr eine Person neu gewählt werden. Wiederwahl ist möglich. 
		\item Der Mitgliederversammlungsvorstand lädt die Mitgliederversammlung ein und stellt die Tagesordnung zusammen mit dem Stammesführer auf. 
		\item Der MVV holt die Berichte der verschiedenen Beauftragten ein und stellt sie der Mitgliederversammlung zur Verfügung. 
		\item Der Vorstand moderiert die Mitgliederversammlung, führt Protokoll und stellt dieses dem Stamm und der Öffentlichkeit innerhalb von vier Wochen zur Verfügung. 
		\item Der MVV wacht über die Beschlüsse der MV und deren Einhaltung. 
	\end{itemize}
    % subsection mitgliederversammlungsvorstand_mvv (end)
    \subsection{Freie Mitarbeiter} % (fold)
    \label{sub:freie_mitarbeiter}
	\begin{itemize}
		\item Jeder Mitarbeiter des Stammes ohne festes Amt ist als freier Mitarbeiter zu sehen. 
		\item Freie Mitarbeiter können bestimmte Aufgaben wahrnehmen, die durch die Führungsrunde beschlossen werden und sich auf Regions- und Landesebene engagieren. 
	\end{itemize}
    % subsection freie_mitarbeiter (end)

% section fuhrungsrunde_und_beauftragungen (end)

\section{Mitgliederversammlung} % (fold)
\label{sec:mitgliederversammlung}

    \subsection{Einberufung und Beschlussfähigkeit} % (fold)
    \label{sub:einberufung_und_beschlussfahigkeit}
	\begin{itemize}
		\item Die Mitgliederversammlung tritt regulär einmal im Jahr zusammen. \sout{und wird durch den Stammesführer einberufen und moderiert. } 
		\item Weitere Mitgliederversammlungen werden auf Verlangen der Führungsrunde oder der einfachen Mehrheit der eingetragenen Mitglieder einberufen. 
		\item Die Mitgliederversammlung setzt sich aus den eingetragenen Mitgliedern des VCP Stammes Martin Luther zusammen, die das zehnte Lebensjahr vollendet haben. 
		\item Die Mitgliederversammlung ist beschlussfähig, wenn sie ordnungsgemäß einberufen wurde. 
		\item Die Einberufung der Mitgliederversammlung erfolgt mindestens 14 Tage vor dem Versammlungstermin durch schriftliche oder elektronische Einladung an jedes Mitglied. 
		\item Wahlen und Beschlüsse erfolgen mit der einfachen Mehrheit ohne Berücksichtigung der Stimmenthaltungen. 
		\item Änderungen, die Stammesordnung betreffend, benötigen eine 2/3 Mehrheit der Anwesenden der Mitgliederversammlung ohne Berücksichtigung der Stimmenthaltungen. 
		\item Die Mitgliederversammlung tagt öffentlich, somit ist es jedem möglich an der Mitgliederversammlung teilzunehmen und Einblick in die Stammesarbeit zu bekommen. 
		\item Jeder ist berechtigt, Anträge an die Mitgliederversammlung zu stellen. 
	\end{itemize}
    % subsection einberufung_und_beschlussfahigkeit (end)
    \subsection{Aufgaben} % (fold)
    \label{sub:aufgaben}
	\begin{itemize}
		\item Die Mitgliederversammlung ist das höchste beschlussfassende Organ des Stammes und entscheidet damit in letzter Instanz innerhalb des Stammes. 
		\item Die Mitgliederversammlung wählt die in dieser Ordnung beschriebenen Beauftragungen. 
		\item Die Mitgliederversammlung nimmt die Berichte der Beauftragen entgegen, der MVV kümmert sich um die Einholung derer. 
		\item Die Mitgliederversammlung kann Entlastungen erteilen. 
	\end{itemize}
    % subsection aufgaben (end)

% section mitgliederversammlung (end)

\section{Die Kassen} % (fold)
\label{sec:die_kassen}

    \subsection{Stammes- und Handkasse} % (fold)
    \label{sub:stammes_und_handkasse}
	\begin{itemize}
		\item Die Stammeskasse dient zur Verwaltung des Stammeskapitals. 
		\item Neben der Stammeskasse kann die Führung einer Handkasse beschlossen werden, welche die Ausgaben aktueller Veranstaltungen decken soll. 
		\item Über die Verwendung der Mittel der Stammes- und Handkasse hat der Kassenwart vollständig Nachweis zu führen. 
		\item Alle Aufgaben in Ausübung eines Amtes im Stamm Martin Luther übernimmt die Stammeskasse. Die Ausgaben sind vor der Führungsrunde zu rechtfertigen. 
		\item Die Ausgaben der Meutenstufe werden durch die Stammeskasse gedeckt. 
	\end{itemize}
    % subsection stammes_und_handkasse (end)

    \subsection{Sippenkassen} % (fold)
    \label{sub:sippenkassen}
	\begin{itemize}
		\item Jeder Sippenführer kann eine Sippenkasse führen. 
		\item Die Kosten von Aktivitäten und Sippenstunden werden durch die Sippenkasse gedeckt. 
		\item Der Sippenführer hat über die Verwendung der Mittel Nachweis zu führen. 
		\item Bei Auflösung der Kasse werden die Mittel den Nachfolgegruppen zur Verfügung gestellt. 
	\end{itemize}
    % subsection sippenkassen (end)

    \subsection{Kassenprüfung} % (fold)
    \label{sub:kassenprufung}
	\begin{itemize}
		\item Die Stammeskasse, die Handkasse und die Sippenkassen sind mindestens einmal jährlich, vor der Mitgliederversammlung, auf die Ordnungsmäßigkeit der Führung hin zu überprüfen. Weitere Kassenprüfungen sind auf Antrag des Stammesführers, der Führungsrunde, oder der einfachen Mehrheit der eingetragenen Mitglieder durchzuführen. 
        \item Bei Auflösung des VCP Stammes Martin Luther Lumdatal sind alle Kassen an den VCP-Hessen, Johannisberg 12, 61231 Bad Nauheim zu überführen.
     \end{itemize}
    % subsection kassenprufung (end)

% section die_kassen (end)
\begin{appendix}
\section{Anhang} % (fold)
\label{sec:anhang}

    \subsection{Versionshistorie} % (fold)
    \label{sub:versionshistorie}
    Die Stammesversammlung wurde auf der ersten Mitgliederversammlung am 10. April 2011 beschlossen. Im folgenden sind die danach getätigten Änderungen beschlossen worden:
        \subsubsection{Mitgliederversammlung am 10. April 2011} % (fold)
        \label{ssub:stammesversammlung_10_april_2011}
        \begin{itemize}
            \item Beschluss der Stammesordnung in der ersten Fassung.
        \end{itemize}
        % subsubsection stammesversammlung_10_april_2011 (end)
    
        \subsubsection{Mitgliederversammlung am 22. Februar 2012} % (fold)
        \label{ssub:stam}
            \begin{itemize}
                \item Ergänzungen: Einjährige Amtszeit der Kassenprüfer, diese war vorher nicht war.
                \item Schaffung der Beauftragung zum Koch. Dieser ist Teil der Stammesleitung und kümmert sich um das leibliche Wohl auf den Veranstaltungen.
            \end{itemize}
        % subsubsection stam (end)
        \subsubsection{Mitgliederversammlung am 31. März 2014} % (fold)
        \label{ssub:mitgliederversammlung_am_31_marz_2014}
            \begin{itemize}
                \item Schaffung einer Möglichkeit zur elektronischen Einladung der Mitgliederversammlung als Alternative zur postalischen Einladung.
                \item Änderung der Zusammensetzung der Mitgliederversammlung: Mitglieder sind erst nach Vollendung des zehnten Lebensjahres (ab elftem Geburtstag) stimmberechtigt.
            \end{itemize}
        % subsubsection mitgliederversammlung_am_31_marz_2014 (end)
    % subsection versionshistorie (end)

% section anhang (end)


\end{appendix}






















\end{document}
